\documentclass[a5paper,10pt]{article}

% https://github.com/InDevRus/filippov-solutions

% Нумерация страниц
\pagenumbering{gobble}

% Настройка полей
\usepackage{geometry}
\geometry{left=1cm}
\geometry{right=1cm}
\geometry{top=1cm}
\geometry{bottom=1.5cm}

% Вставка таблицы
\usepackage{array}
\usepackage{tabularx}

% Инструменты выравнивания формул
\usepackage{amsmath}
\usepackage{mathtools}

% Диакритика в математике
\usepackage{amsfonts}

% Вычисления размеров на ходу
\usepackage{calc}

% Удобные записи производных
\usepackage[italic=true]{derivative}

% Настройки сносок
\usepackage{enotez}

% Длинные знаки векторов
\usepackage{anyfontsize}
\usepackage[b]{esvect}

% Вставка картинок
\usepackage{graphicx}

% Вставка красной строки
\usepackage{indentfirst}
\setlength{\parindent}{1em}

% Условия в командах
\usepackage{xifthen}

% Луа-скрипты
\usepackage{luacode}

% Настройка команд отдельно в математическом и текстовом режимах
\usepackage{mathcommand}

% Для повторения LaTeX кода
\usepackage{multido}

% Конфигурация отступов формул
\delimitershortfall-1sp
\usepackage{mleftright}
\mleftright

% Растягивание объектов
\usepackage{relsize}

% Обтекание текстом
\usepackage{wrapfig}
\usepackage{subcaption}
\usepackage{floatflt}

% Поддержка ввода кириллицы
\usepackage[english, russian]{babel}

% Настройка корневой позиции для компилятора
\usepackage{import}
\providecommand*{\ProjectRootPath}{..}

\import{\ProjectRootPath/preambles/macros/}{theorems}

\import{\ProjectRootPath/preambles/fonts}{font_setup}

% Знак градуса
\usepackage{siunitx}

\import{\ProjectRootPath/preambles/macros/}{operators}
\import{\ProjectRootPath/preambles/macros/}{redefinitions}

\import{\ProjectRootPath/preambles/fonts}{letters_setup}
\import{\ProjectRootPath/preambles/fonts/adjustments}{custom_superscripts}

\import{\ProjectRootPath/preambles/custom}{formatting}
\import{\ProjectRootPath/preambles/custom}{frames}
\import{\ProjectRootPath/preambles/custom}{list_setup}

% TODO: перевести команды на современный стиль объявления

\begin{document}
    $$\tsys{& \parder{z}{x} = 2yz - \pow{z}{2} \\& \parder{z}{y} = xz}$$

    $x \parder{z}{x} = 2y \cdot xz - x\pow{z}{2} = 2y \parder{z}{y} - x\pow{z}{2}$.
    $x \parder{z}{x} - 2y \parder{z}{y} = -x\pow{z}{2}$.
    $\dfrac {dx} {x} = \dfrac {dy} {-2y} = \dfrac {dz} {-x\pow{z}{2}}$.

    В частности, интегрируемыми комбинациями являются $\dfrac {dx} {x} = \dfrac {dy} {-2y}$ и $\dfrac {dx} {x} = \dfrac {dz} {-x\pow{z}{2}}$.
    Первая даёт $2\ln\vbr{x} + \ln\vbr{y} = A$, $\pow{x}{2} y = A$; вторая даёт $dx + \dfrac {dz} {\pow{z}{2}} = 0$,
    $x - \dfrac {1} {z} = B$. При делении было также потеряно решение $z\br{x; y} = 0$. \vspace{1mm}

    Таким образом, общее решение уравнения
    $x \parder{z}{x} - 2y \parder{z}{y} = -x\pow{z}{2}$
    имеет вид
    $x - \dfrac {1} {z} = \linebreak = f\br{\pow{x}{2} y}$, то есть
    $z\br{x; y} = \dfrac {1} {x + f\br{\pow{x}{2} y}}$, а также $z\br{x; y} = 0$.

    Покажем, что никакое решение вида $z\br{x; y} = \dfrac {1} {x + f\br{\pow{x}{2} y}}$ не удовлетворяет условию задачи.

    $\parder{z}{x}
    = \dfrac {\partial} {\partial x} \br{\dfrac {1} {x + f\br{\pow{x}{2} y}}}
    = \dfrac {-1 - 2x y \f\'\br{\pow{x}{2} y}} {\br{x + f\br{\pow{x}{2} y}}^2} = 2yz - \pow{z}{2}$.

    $\parder{z}{y} = \dfrac {-\pow{x}{2} f\'\br{\pow{x}{2} y}} {\br{x + f\br{\pow{x}{2} y}}^2} = x z$.
    $\dfrac {-x \f\'\br{\pow{x}{2} y}} {\br{x + f\br{\pow{x}{2} y}}^2} = z$.

    Поделив последние равенства, получим
    $\dfrac {1 + 2xy f\'\br{\pow{x}{2}y}} {x f\'\br{\pow{x}{2}y}} = 2y - z$. $1 = -z x f\'\br{\pow{x}{2}y}$.
    $x f\'\br{\pow{x}{2}y} = -x - f\br{\pow{x}{2} y}$.

    Итак, для некоторого дифференцируемого выражения $f\br{t}$ должно выполняться \linebreak
    $x \br{ \f\'\br{\pow{x}{2} y} + 1} = -f\br{\pow{x}{2} y}$ для всех $x$ и $y$.

    Выражения $f\'\br{\pow{x}{2} y} + 1$ и $f\br{\pow{x}{2} y}$ не могут одновременно тождественно равняться нулю.
    Следовательно, вдоль какой-либо кривой вида $\pow{x}{2} y = \tau$, последнее равенство невозможно, так как такая фиксация позволяет однозначно выразить $x$ из равенства
    $x \br{\f\'\br{\tau} + 1} = -f\br{\tau}$,
    то есть равенство выполняется не более, чем в одной точке. А поскольку на этих кривых тождество невозможно, оно тем более не может выполняться и в более общем случае.

    Итак, единственным допустимым решением является $z\br{x; y} = 0$.

\end{document}
